\begin{lemma}
\label{editorial-source-lemma-Galois-orbit}
Let $K/k$ be an extension of fields.
Let $\overline{k}/k$ be a separable algebraic closure.
Then $\text{Gal}(\overline{k}/k)$ acts transitively on the
primes of $\overline{k} \otimes_k K$.
\end{lemma}

\begin{proof}
Retain the original separable closure
$\overline k$, put $G=\operatorname{Gal}(\overline k/k)$,
and let $K/k'/k$ be the original
subextension of Lemma
\ref{editorial-source-lemma-make-geometrically-irreducible}.
Write
$$
B=\overline k\otimes_k k',\qquad
A=\overline k\otimes_k K.
$$
Both tensors are nonzero by extension
of vector-space bases over $k$.
As $\overline k/k$ is algebraic,
these are integral algebras over
the respective fields $k'$ and
$K$: each generator $a\otimes1$
satisfies its original monic
polynomial over $k$. Thus all
their primes are maximal by
Lemma \ref{editorial-source-lemma-integral-no-inclusion},
and are also minimal since
there are no strict inclusions
between primes over a field.
The two original comparison maps
$$
B\otimes_{k'}K\longrightarrow A,
\qquad(a\otimes b)\otimes c\longmapsto a\otimes bc,
$$
and $a\otimes c\mapsto(a\otimes1)\otimes c$
are inverse by $k'$-balancing.
Since $K/k'$ is geometrically
irreducible, Lemma
\ref{editorial-source-lemma-geometrically-irreducible-any-base-change}
therefore identifies the spectra
of $A$ and $B$, including their
topologies by its minimal-prime
homeomorphism. The map is
contraction, and commutes with
the original automorphisms
$g\otimes1$ for $g\in G$.

For a prime $\mathfrak p$ of $B$,
its residue field is separably
algebraic over the original
$\overline k$. Indeed $B$ is
generated over $\overline k$ by
$1\otimes b$ for $b\in k'$,
and each of their images satisfies
the original separable polynomial
of $b$ over $k$. Minimal
polynomials over $\overline k$
divide these separable polynomials.
Every element of the field
$B/\mathfrak p$ belongs to a
finite subextension generated
by finitely many such images,
so that extension is separable.
As $\overline k$ is separably
closed, its structural embedding
into $B/\mathfrak p$ is an
isomorphism. Its inverse composed
with $b\mapsto[1\otimes b]$
defines a unique embedding
$\sigma:k'\to\overline k$ over $k$.
Consequently
$$
\mathfrak p=\mathfrak p_\sigma
=\ker(\mu_\sigma),\qquad
\mu_\sigma:B\to\overline k,\quad
\sum_i a_i\otimes b_i\longmapsto\sum_i a_i\sigma(b_i).
$$
Conversely this map is surjective
because $a=\mu_\sigma(a\otimes1)$,
so its kernel is a prime. The
structural field isomorphism
above proves uniqueness of
$\sigma$ from its kernel.
This gives the stated bijection
with $\operatorname{Hom}_k(k',\overline k)$.
Separability, rather than
maximality alone, is what
identifies these residue fields
with the separably closed field.

Given embeddings $\sigma,\tau$,
extend the isomorphism
$\tau\sigma^{-1}:\sigma(k')\to\tau(k')$
to a $k$-embedding of $\overline k$
in itself. Here are the needed
extension details, including
infinite $k'/k$. Order partial
extensions by inclusion and
take unions on chains. If a
maximal partial extension has
domain $L\neq\overline k$, choose
$a\in\overline k\setminus L$.
Its minimal polynomial over
$L$ divides a separable
polynomial over $k$, hence is
separable. Transporting its
coefficients gives a separable
polynomial over the image field,
which has a root in the
separably closed $\overline k$.
Sending $a$ to that root extends
the embedding to $L(a)$, a
contradiction. Thus the domain
is all of $\overline k$.
Any $k$-embedding of this field
in itself is surjective: each
$b\in\overline k$ lies in a
finite Galois subextension $N/k$,
obtained by adjoining all roots
of its separable minimal polynomial.
The embedding permutes those
roots, sends $N$ into itself,
and its injective $k$-linear
map on the finite-dimensional
$N$ is surjective. In particular
$b$ is in its image. We have
therefore constructed $g\in G$
with $g\sigma=\tau$.

Use the action on ideals
$\mathfrak p\mapsto(g\otimes1)(\mathfrak p)$.
For every original $z\in B$,
$$
\mu_{g\sigma}((g\otimes1)z)=g(\mu_\sigma(z)).
$$
It follows that
$(g\otimes1)(\mathfrak p_\sigma)=\mathfrak p_{g\sigma}$.
This proves transitivity on
$\Spec(B)$ and, by equivariant
contraction, on the original
$\Spec(A)$.

The correspondence also identifies
the actual prime and residue
field of $A$. Give $\overline k$
its $k'$-structure through
$\sigma$. The extension
$\overline k/\sigma(k')$ is
separable algebraic. The tensor-field
form of Lemma
\ref{editorial-source-lemma-geometrically-irreducible-separable-elements}
applied to $K/k'$ therefore makes
$$
C_\sigma=\overline k\otimes_{k',\sigma}K
$$
a field. The original quotient
comparison sends the class of
$a\otimes x$ to $a\otimes x$
and gives an isomorphism
$A/\mathfrak p_\sigma A\simeq C_\sigma$;
its inverse sends the same pure
tensor to its class. The additional
$k'$-balancing relations are
exactly the extended kernel
of $\mu_\sigma$, so both maps
are well-defined. Hence the
corresponding prime is exactly
$\mathfrak q_\sigma=\mathfrak p_\sigma A$,
and its residue field is
$C_\sigma$. No further radical
or residue-field identification
is needed or assumed.

For a fixed $\sigma$, its
stabilizer in $G$ is exactly
the subgroup fixing $\sigma(k')$
pointwise, by the injectivity
of the embedding-to-prime
correspondence. Thus the orbit
is the corresponding coset
set. If $K/k$ is finitely
generated, the preceding lemma
makes $k'/k$ finite separable;
Fields, Lemma
\ref{fields-lemma-separable-equality}
then gives exactly $[k':k]$
primes of $A$. Since this
finite spectrum has every
point closed, each point is
also open: its complement is
a finite union of closed
points. It is therefore a
finite discrete orbit. No
finite-degree count is asserted
for an infinite $k'/k$.
\end{proof}